\subsection{Localization of associated prime ideals}

PROPOSITION 5. Let A be a ring, S a multiplicative subset of A, \(\Phi\) the set of prime ideals of A which do not meet S and M an A-module. Then:

\begin{enumerate}
\def\labelenumi{(\roman{enumi})}
\item
  The mapping \(\mathfrak{p} \mapsto \mathrm{S}^{-1}\mathfrak{p}\) is a bijection of \(\operatorname{Ass}_{\mathbf{A}}(\mathbf{M}) \cap \Phi\) onto a subset of \(\operatorname{Ass}_{\mathbf{S}^{-1}\mathbf{A}}(\mathbf{S}^{-1}\mathbf{M})\) .
\item
  If \(\mathfrak{p} \in \Phi\) is a finitely generated ideal and \(S^{-1}\mathfrak{p} \in \operatorname{Ass}_{\mathbf{S}^{-1}\mathbf{A}}(\mathbf{S}^{-1}\mathbf{M})\) , then \(\mathfrak{p} \in \operatorname{Ass}_A(M)\) .
\end{enumerate}

Recall (Chapter II, § 2, no. 5, Proposition 11) that the mapping \(\mathfrak{p} \mapsto \mathrm{S}^{-1}\mathfrak{p}\) is a bijection of \(\Phi\) onto the set of prime ideals of \(\mathrm{S}^{-1}\mathrm{A}\) . If \(\mathfrak{p} \in \operatorname{Ass}_{\mathbf{A}}(\mathbf{M}) \cap \Phi\) , \(\mathfrak{p}\) is the annihilator of a monogenous submodule \(\mathbf{N}\) of \(\mathbf{M}\) ; then \(\mathrm{S}^{-1}\mathfrak{p}\) is the annihilator of the monogenous submodule \(\mathrm{S}^{-1}\mathrm{N}\) of \(\mathrm{S}^{-1}\mathrm{M}\) (Chapter II, § 2, no. 4, formula (9)) and hence \(\mathbf{S}^{-1}\mathfrak{p}\in\mathrm{Ass}_{\mathbf{S}^{-1}\mathbf{A}}(\mathbf{S}^{-1}\mathbf{M})\) . Conversely, suppose that \(\mathfrak{p}\in\Phi\) is finitely generated and such that \(\mathbf{S}^{-1}\mathfrak{p}\) is associated with \(\mathbf{S}^{-1}\mathbf{M}\) ; then there exists \(x\in\mathbf{M}\) and \(t\in\mathbf{S}\) such that \(\mathbf{S}^{-1}\mathfrak{p}\) is the annihilator of \(x/t\) . Let \((a_{i})_{1\leqslant i\leqslant n}\) be a system of generators of \(\mathfrak{p}\) ; then \((a_{i}/1)(x/t)=0\) and hence there exists \(s_{i}\in\mathbf{S}\) such that \(s_{i}a_{i}x=0\) ( \(1\leqslant i\leqslant n\) ). Let us write \(s=s_{1}s_{2}\ldots s_{n}\) ; for all \(a\in\mathfrak{p}\) , \(sax=0\) , whence \(\mathfrak{p}\subset\mathrm{Ann}(sx)\) ; on the other hand, if \(b\in\mathbf{A}\) satisfies \(bsx=0\) , then \(b/1\in\mathbf{S}^{-1}\mathfrak{p}\) by definition, whence \(b\in\mathfrak{p}\) . Then \(\mathfrak{p}=\mathrm{Ann}(sx)\) and \(\mathfrak{p}\in\mathrm{Ass}_{\mathbf{A}}(\mathbf{M})\) .

COROLLARY. If the ring A is Noetherian, the mapping \(\mathfrak{p} \mapsto \mathrm{S}^{-1}\mathfrak{p}\) is a bijection of \(\operatorname{Ass}_{\mathbf{A}}(\mathbf{M}) \cap \Phi\) onto \(\operatorname{Ass}_{\mathbf{S}^{-1}\mathbf{A}}(\mathbf{S}^{-1}\mathbf{M})\) .

If A is not Noetherian, the mapping \(p \mapsto S^{-1}p\) of \(\operatorname{Ass}_{\mathbf{A}}(\mathbf{M}) \cap \Phi\) to \(\operatorname{Ass}_{\mathbf{S}^{-1}\mathbf{A}}(\mathbf{S}^{-1}\mathbf{M})\) is not necessarily surjective (Exercise 1).

PROPOSITION 6. Let A be a Noetherian ring, M an A-module, S a multiplicative subset of A and \(\Psi\) the set of elements of \(\operatorname{Ass}_{\mathbf{A}}(\mathbf{M})\) which do not meet S. Then the kernel N of the canonical mapping \(M \to S^{-1}M\) is the unique submodule of M which satisfies the relations

\[
\operatorname{Ass} (N) = \operatorname{Ass} (M) - \Psi , \quad \operatorname{Ass} (M / N) = \Psi .\tag{3}
\]

By Proposition 4 of no. 1, there exists a submodule \(\mathbf{N}'\) of \(\mathbf{M}\) which satisfies the relations \(\operatorname{Ass}(\mathbf{N}') = \operatorname{Ass}(\mathbf{M}) - \Psi'\) and \(\operatorname{Ass}(\mathbf{M}/\mathbf{N}') = \Psi'\) . We need to prove \(\mathbf{N}' = \mathbf{N}\) . Consider the commutative diagram

\[
\begin{array}{c} \mathrm{M} \xrightarrow {p} \mathrm{M} / \mathrm{N} ^ {\prime} \\ u \Biggl \downarrow \\ \mathrm{S} ^ {- 1} \mathrm{M} \xrightarrow [ \mathrm{S} ^ {- 1 p} ]{} \mathrm{S} ^ {- 1} (\mathrm{M} / \mathrm{N} ^ {\prime}) \end{array}
\]

where \(p, u, v\) are the canonical homomorphisms. We shall show that \(S^{-1}p\) and \(v\) are injective, which will prove that \(u\) and \(p\) have the same kernel and hence \(N' = N\) .

As \(\operatorname{Ass}(\mathbf{N}') \cap \Psi' = \varnothing\) , every element of \(\operatorname{Ass}(\mathbf{N}')\) meets S. Then \(\operatorname{Ass}_{\mathbf{S}^{-1}\mathbf{A}}(\mathbf{S}^{-1}\mathbf{N}') = \varnothing\) (Corollary to Proposition 5), whence \(\mathbf{S}^{-1}\mathbf{N}' = \{0\}\) (no. 1, Corollary 1 to Proposition 2), which proves that \(\mathbf{S}^{-1}\mathfrak{p}\) is injective (Chapter II, § 2, no. 4, Theorem 1). On the other hand, if \(x\) belongs to the kernel K of \(v\) , then \(\operatorname{Ann}(x) \cap \mathbf{S} \neq \varnothing\) (Chapter II, § 2, no. 2, Proposition 4); hence \(\operatorname{Ass}(\mathbf{K}) = \varnothing\) since \(\operatorname{Ass}(\mathbf{K}) \subset \operatorname{Ass}(\mathbf{M}/\mathbf{N}') = \Psi'\) ; we deduce that \(\mathbf{K} = \{0\}\) (no. 1, Corollary 1 to Proposition 2) and \(v\) is injective.

